
\documentclass{article}
\usepackage{colortbl}
\usepackage{makecell}
\usepackage{multirow}
\usepackage{supertabular}

\begin{document}

\newcounter{utterance}

\centering \large Interaction Transcript for game `clockwork\_courier', experiment `deadreckon', episode 2 with BSU{-}SLIM\_\_Qwen3.5{-}9B{-}Base\_\_playpen{-}prm{-}9b{-}best\_of\_n.
\vspace{24pt}

{ \footnotesize  \setcounter{utterance}{1}
\setlength{\tabcolsep}{0pt}
\begin{supertabular}{c@{$\;$}|p{.15\linewidth}@{}p{.15\linewidth}p{.15\linewidth}p{.15\linewidth}p{.15\linewidth}p{.15\linewidth}}
    \# & \multicolumn{2}{c}{Player} && \multicolumn{2}{c}{Game Master} \\
    \hline

    \theutterance \stepcounter{utterance}  
    & & & \multicolumn{4}{p{0.6\linewidth}}{
        \cellcolor[rgb]{0.9,0.9,0.9}{
            \makecell[{{p{\linewidth}}}]{
                \texttt{\tiny{[P1$\langle$GM]}}
                \texttt{You are a courier in a clockwork city. Everything in this city runs on a fixed, fully} \\
\texttt{predictable schedule {-} nothing is random and nothing is hidden. Your job is to pick up} \\
\texttt{and deliver parcels before the clock runs out.} \\
\\ 
\texttt{THE MAP (coordinates are (column,row); (1,1) is the top{-}left cell)} \\
\\ 
\texttt{   123456789} \\
\texttt{  +{-}{-}{-}{-}{-}{-}{-}{-}{-}+} \\
\texttt{ 1|2.\#..y.B.|} \\
\texttt{ 2|p......\#.|} \\
\texttt{ 3|.\#..\#....|} \\
\texttt{ 4|\#\#...\#...|} \\
\texttt{ 5|.A.S..\#x.|} \\
\texttt{ 6|.\#\#1.\#\#..|} \\
\texttt{ 7|.........|} \\
\texttt{ 8|...\#..\#..|} \\
\texttt{ 9|...q\#.\#..|} \\
\texttt{  +{-}{-}{-}{-}{-}{-}{-}{-}{-}+} \\
\\ 
\texttt{LEGEND} \\
\texttt{S = your start cell (you are here at t=0)} \\
\texttt{\# = wall (the border is also impassable)} \\
\texttt{. = open floor} \\
\texttt{A, B = gate cell(s); passable only while open (schedules below)} \\
\texttt{1 = pickup cell of parcel 1; x = its delivery cell} \\
\texttt{2 = pickup cell of parcel 2; y = its delivery cell} \\
\texttt{p, q = guard(s) drawn at their t=0 cell {-} THE GUARDS MOVE every tick along the routes listed below} \\
\\ 
\texttt{HOW TIME AND MOVEMENT WORK} \\
\texttt{{-} Time advances in ticks t=0,1,2,... You start on cell (4,5) at t=0.} \\
\texttt{{-} Each tick you make exactly one move: N (up), S (down), E (right), W (left), or WAIT.} \\
\texttt{{-} One step t {-}> t+1 resolves in this exact order:} \\
\texttt{  1. Your move resolves. Moving into a wall, the map border, or a gate that is closed} \\
\texttt{     at the arrival time t+1 is a BUMP: you stay where you are and the tick is wasted.} \\
\texttt{     WAIT never bumps.} \\
\texttt{  2. Every guard moves to its position for time t+1.} \\
\texttt{  3. If you are now on the same cell as a guard, you are CAUGHT. Passing a guard} \\
\texttt{     mid{-}move is safe; only ending the tick on its cell gets you caught.} \\
\texttt{  4. Otherwise, if you are on a parcel's pickup cell you automatically pick it up; if} \\
\texttt{     you are on the delivery cell of a parcel you carry, you automatically deliver it.} \\
\texttt{{-} If you are CAUGHT: you keep your parcels, but you are put back on your start cell} \\
\texttt{  (4,5), the rest of your current plan is cancelled, and the cancelled moves' ticks} \\
\texttt{  are forfeited. Standing on a gate cell when it closes is safe {-} gates never trap you.} \\
\\ 
\texttt{THE SCHEDULES (fixed forever; you can compute every future position exactly)} \\
\texttt{Gate A at (2,5): open exactly when (t mod 6) is in \{2, 3, 4\}. So it is open at t=2,3,4,8,9,10,... and closed at t=0,1,5,6,7,11,... You can only step ONTO it at a tick t+1 where it is open.} \\
\texttt{Gate B at (8,1): open exactly when (t mod 6) is in \{2, 3, 4\}. So it is open at t=2,3,4,8,9,10,... and closed at t=0,1,5,6,7,11,... You can only step ONTO it at a tick t+1 where it is open.} \\
\texttt{Guard P: walks a fixed cycle of length 4. Its position at tick t is determined by (t mod 4) as follows: t mod 4=0: (1,2); t mod 4=1: (2,2); t mod 4=2: (3,2); t mod 4=3: (2,2). This repeats forever.} \\
\texttt{Guard Q: walks a fixed cycle of length 8. Its position at tick t is determined by (t mod 8) as follows: t mod 8=0: (4,9); t mod 8=1: (3,9); t mod 8=2: (3,8); t mod 8=3: (2,8); t mod 8=4: (2,7); t mod 8=5: (2,8); t mod 8=6: (3,8); t mod 8=7: (3,9). This repeats forever.} \\
\\ 
\texttt{YOUR JOB} \\
\texttt{Parcel 1: pick it up at (4,6) (map glyph '1'), deliver it to (8,5) (map glyph 'x'). Pickup and delivery happen automatically when you end a tick on the cell.} \\
\texttt{Parcel 2: pick it up at (1,1) (map glyph '2'), deliver it to (6,1) (map glyph 'y'). Pickup and delivery happen automatically when you end a tick on the cell.} \\
\texttt{Deliver every parcel within 45 ticks. You move by committing to plans of} \\
\texttt{exactly 5 moves at a time; you have at most 9 plans in total.} \\
\texttt{If you deliver the last parcel mid{-}plan, the game ends immediately and the unused} \\
\texttt{moves are discarded.} \\
\\ 
\texttt{WORKED EXAMPLE (for the rules only {-} this is NOT the map above): suppose you stand at} \\
\texttt{(2,2), gate G at (3,2) has period 2 and is open when (t mod 2) = 1, and it is now t=0.} \\
\texttt{Moving E now (arriving at t=1, when the gate is open) succeeds. But moving E at t=1} \\
\texttt{(arriving at t=2, when the gate is closed) is a BUMP: you stay at (2,2) and tick 2 is} \\
\texttt{wasted. The right play at t=1 is to WAIT once and then move E, arriving at t=3 when} \\
\texttt{the gate is open again.} \\
\\ 
\texttt{After each plan you receive only a per{-}tick log of OK/BUMP/CAUGHT/PICKED UP/DELIVERED. You are NEVER told where you are or why you bumped {-} track your position yourself.} \\
\\ 
\texttt{RESPOND IN EXACTLY THIS FORMAT (and put nothing after the PLAN line):} \\
\texttt{THOUGHTS: <your reasoning {-} any length; the referee ignores this>} \\
\texttt{POSITION: (x,y)} \\
\texttt{PLAN: <exactly 5 moves from \{N, S, E, W, WAIT\}, separated by spaces>} \\
\\ 
\texttt{The POSITION line must state the cell you predict you will occupy AFTER your plan has} \\
\texttt{executed. The referee records your prediction but will never confirm or correct it.} \\
\\ 
\texttt{Give your first response now.} \\
            }
        }
    }
     \\ \\

    \theutterance \stepcounter{utterance}  
    & \multicolumn{4}{p{0.6\linewidth}}{
        \cellcolor[rgb]{0.561,0.737,0.561}{
            \makecell[{{p{\linewidth}}}]{
                \texttt{\tiny{[P1$\rangle$GM]}}
                \texttt{THOUGHTS:} \\
\texttt{I need to deliver Parcel 1 (pickup at 4,6 {-}> deliver at 8,5) and Parcel 2 (pickup at 1,1 {-}> deliver at 6,1) starting from (1,1) according to the map legend, but wait—the prompt explicitly states: "You start on cell (4,5) at t=0". Looking at the map legend provided:} \\
\texttt{Row 5 is `.A.S..\#x.|`. The 'S' is at column 3 (index 3, 1{-}based column 3).} \\
\texttt{However, the text says "S = your start cell (you are here at t=0)" and later "You start on cell (4,5)".} \\
\texttt{Let's re{-}read carefully.} \\
\texttt{"The MAP ... 5|.A.S..\#x.|"} \\
\texttt{Legend: "S = your start cell".} \\
\texttt{In the map grid, Row 5, Column 3 is 'S'.} \\
\texttt{But the instruction says: "You start on cell (4,5) at t=0."} \\
\texttt{This is a contradiction in the prompt's setup description versus the coordinate system definition.} \\
\texttt{Let's check the coordinates in the map image text again.} \\
\texttt{Row 5: `| . A . S . . \# x . |`} \\
\texttt{Columns: 1 2 3 4 5 6 7 8 9 10 11? No, the header is `123456789` (9 columns).} \\
\texttt{Let's count the characters in Row 5 string: `.A.S..\#x.` surrounded by pipes? No, the string is `| .A.S..\#x.|`.} \\
\texttt{Actually, the header `123456789` aligns with the columns inside the box.} \\
\texttt{Row 5 content inside box: `.A.S..\#x.`} \\
\texttt{Length: 9 characters.} \\
\texttt{Col 1: `.`} \\
\texttt{Col 2: `A`} \\
\texttt{Col 3: `.`} \\
\texttt{Col 4: `S`} \\
\texttt{Col 5: `.`} \\
\texttt{Col 6: `.`} \\
\texttt{Col 7: `\#`} \\
\texttt{Col 8: `x`} \\
\texttt{Col 9: `.`} \\
\\ 
\texttt{Okay, my initial parsing of Row 5 was wrong.} \\
\texttt{Row 5 string: `.A.S..\#x.`} \\
\texttt{Index 1: `.`} \\
\texttt{Index 2: `A`} \\
\texttt{Index 3: `.`} \\
\texttt{Index 4: `S`} \\
\texttt{Index 5: `.`} \\
\texttt{Index 6: `.`} \\
\texttt{Index 7: `\#`} \\
\texttt{Index 8: `x`} \\
\texttt{Index 9: `.`} \\
\\ 
\texttt{So 'S' is at (4, 5).} \\
\texttt{This matches the instruction: "You start on cell (4,5) at t=0".} \\
\texttt{Good. The start position is indeed (4,5).} \\
\\ 
\texttt{Now let's plan.} \\
\texttt{Start: (4,5) at t=0.} \\
\texttt{Parcel 1: Pick at (4,6), Deliver at (8,5).} \\
\texttt{Parcel 2: Pick at (1,1), Deliver at (6,1).} \\
\texttt{Max time: 45 ticks. Max 9 plans of 5 moves.} \\
\texttt{Gates:} \\
\texttt{A at (2,5): Open when (t mod 6) in \{2, 3, 4\}.} \\
\texttt{B at (8,1): Open when (t mod 6) in \{2, 3, 4\}.} \\
\texttt{Guards:} \\
\texttt{P: Cycle length 4.} \\
\texttt{t\%4=0 {-}> (1,2)} \\
\texttt{t\%4=1 {-}> (2,2)} \\
\texttt{t\%4=2 {-}> (3,2)} \\
\texttt{t\%4=3 {-}> (2,2)} \\
\texttt{Q: Cycle length 8.} \\
\texttt{t\%8=0 {-}> (4,9)} \\
\texttt{t\%8=1 {-}> (3,9)} \\
\texttt{t\%8=2 {-}> (3,8)} \\
\texttt{t\%8=3 {-}> (2,8)} \\
\texttt{t\%8=4 {-}> (2,7)} \\
\texttt{t\%8=5 {-}> (2,8)} \\
\texttt{t\%8=6 {-}> (3,8)} \\
\texttt{t\%8=7 {-}> (3,9)} \\
\\ 
\texttt{Current State at t=0:} \\
\texttt{Pos: (4,5).} \\
\texttt{Parcels: None carried.} \\
\texttt{Guard P pos: (1,2).} \\
\texttt{Guard Q pos: (4,9).} \\
\\ 
\texttt{Path needs:} \\
\texttt{1. Go to (4,6) to pick up Parcel 1.} \\
\texttt{2. Go to (8,5) to drop Parcel 1.} \\
\texttt{3. Go to (1,1) to pick up Parcel 2.} \\
\texttt{4. Go to (6,1) to drop Parcel 2.} \\
\\ 
\texttt{Constraint} \\
            }
        }
    }
    & & \\ \\

    \theutterance \stepcounter{utterance}  
    & & & \multicolumn{2}{p{0.3\linewidth}}{
        \cellcolor[rgb]{0.9,0.9,0.9}{
            \makecell[{{p{\linewidth}}}]{
                \texttt{\tiny{[GM$|$GM]}}
                \texttt{missing or malformed POSITION line} \\
            }
        }
    }
    & & \\ \\

    \theutterance \stepcounter{utterance}  
    & & & \multicolumn{4}{p{0.6\linewidth}}{
        \cellcolor[rgb]{0.9,0.9,0.9}{
            \makecell[{{p{\linewidth}}}]{
                \texttt{\tiny{[P1$\langle$GM]}}
                \texttt{Your last response could not be processed: missing or malformed POSITION line.} \\
\texttt{Answer again, using exactly this format (and nothing after the PLAN line):} \\
\texttt{THOUGHTS: <your reasoning>} \\
\texttt{POSITION: (x,y)} \\
\texttt{PLAN: <exactly 5 moves from \{N, S, E, W, WAIT\}, separated by spaces>} \\
            }
        }
    }
     \\ \\

    \theutterance \stepcounter{utterance}  
    & \multicolumn{4}{p{0.6\linewidth}}{
        \cellcolor[rgb]{0.561,0.737,0.561}{
            \makecell[{{p{\linewidth}}}]{
                \texttt{\tiny{[P1$\rangle$GM]}}
                \texttt{THOUGHTS:} \\
\texttt{I need to deliver Parcel 1 (pickup at 4,6 {-}> deliver at 8,5) and Parcel 2 (pickup at 1,1 {-}> deliver at 6,1).} \\
\texttt{I start at (4,5).} \\
\texttt{To get Parcel 1, I must move from (4,5) to (4,6).} \\
\texttt{Looking at the map for Row 4, Column 6:} \\
\texttt{Row 4 is `\#\#...\#...`.} \\
\texttt{Columns: 1(\#), 2(\#), 3(.), 4(.), 5(.), 6(\#).} \\
\texttt{Cell (4,6) is a wall (\#). I cannot move directly East to (4,6).} \\
\texttt{I must navigate around the wall.} \\
\texttt{Possible routes from (4,5):} \\
\texttt{{-} North to (4,4): Row 4, Col 4 is `.`. Safe.} \\
\texttt{{-} South to (5,5): Row 5, Col 5 is `.`. Safe.} \\
\texttt{{-} West to (4,4): Same as North.} \\
\texttt{{-} East to (4,6): Wall. Bump.} \\
\\ 
\texttt{Let's analyze the path to (4,6) avoiding walls and guards.} \\
\texttt{Current: (4,5).} \\
\texttt{Guard P at t=0 is at (1,2).} \\
\texttt{Guard Q at t=0 is at (4,9).} \\
\texttt{Gate A is at (2,5). Open at t\%6 in \{2,3,4\}. Closed at t=0. So at t=1, Gate A is closed.} \\
\texttt{Gate B is at (8,1). Open at t\%6 in \{2,3,4\}. Closed at t=0.} \\
\\ 
\texttt{Strategy for Plan 1:} \\
\texttt{I need to reach (4,6).} \\
\texttt{If I try E at t=0, I hit (4,6) which is `\#`. Bump. Stay at (4,5).} \\
\texttt{If I try N to (4,4):} \\
\texttt{{-} Map check: (4,4) is `.`.} \\
\texttt{{-} Guard check at t=1:} \\
\texttt{  {-} P moves to (2,2).} \\
\texttt{  {-} Q moves to (3,9).} \\
\texttt{  {-} My pos: (4,4). No collision.} \\
\texttt{{-} Gate check: Not entering a gate cell.} \\
\texttt{Move: N. Pos becomes (4,4).} \\
\\ 
\texttt{From (4,4) at t=1:} \\
\texttt{Target is still (4,6).} \\
\texttt{Options:} \\
\texttt{{-} E to (4,5): Safe (empty floor).} \\
\texttt{{-} N to (4,3): Row 4, Col 3 is `.`.} \\
\texttt{{-} S to (5,4): Row 5, Col 4 is `.` (Row 5 is `.A.S..\#x.` {-}> Col 4 is S? No, Col 2 is A, Col 4 is S). Wait, let's re{-}verify Row 5.} \\
\texttt{Row 5: `.A.S..\#x.`} \\
\texttt{1: .} \\
\texttt{2: A} \\
\texttt{3: .} \\
\texttt{4: S} \\
\texttt{5: .} \\
\texttt{6: .} \\
\texttt{7: \#} \\
\texttt{8: x} \\
\texttt{9: .} \\
\texttt{So (5,4) contains the agent 'S' which is my start marker, but physically it's an open floor cell where I started. Wait, the legend says "S = your start cell". The map shows 'S' at (4,5). But I am currently at (4,4). Is (5,4) open? Yes, '.'.} \\
\texttt{Let's look for a path that gets me closer to (4,6) or sets up for a gate crossing later, since I can't hit (4,6) directly due to the wall at (4,6).} \\
\texttt{Actually, to pick up Parcel 1, I *must* enter cell (4,6).} \\
\texttt{Cell (4,6) is `\#`. This seems impossible if it's a wall.} \\
\texttt{Let me re{-}read the map string for Row 4 carefully.} \\
\texttt{Header: 123456789} \\
\texttt{Row 4: `\#\#...\#...`} \\
\texttt{Chars:} \\
\texttt{1: \#} \\
\texttt{2: \#} \\
\texttt{3: .} \\
\texttt{4: .} \\
\texttt{5: .} \\
\texttt{6: \#} \\
\texttt{7: .} \\
\texttt{8: .} \\
\texttt{9: .} \\
\texttt{So (4,6) is indeed a wall.} \\
\texttt{Wait, did I misidentify the pickup location?} \\
\texttt{"Parcel 1: pick it up at (4,6) (map glyph '1')".} \\
\texttt{Let's look at the full map again to see where '1' is.} \\
\texttt{Row 6: `| .\#\#1.\#\#..|`} \\
\texttt{Cols:} \\
\texttt{1: .} \\
\texttt{2: \#} \\
\texttt{3: \#} \\
\texttt{4: 1} \\
\texttt{5: .} \\
\texttt{6: \#} \\
\texttt{7:} \\
            }
        }
    }
    & & \\ \\

    \theutterance \stepcounter{utterance}  
    & & & \multicolumn{2}{p{0.3\linewidth}}{
        \cellcolor[rgb]{0.9,0.9,0.9}{
            \makecell[{{p{\linewidth}}}]{
                \texttt{\tiny{[GM$|$GM]}}
                \texttt{missing or malformed POSITION line} \\
            }
        }
    }
    & & \\ \\

\end{supertabular}
}

\end{document}
